\documentclass[11pt, a4paper]{article}

\usepackage[a4paper, top=2.5cm, bottom=2.5cm, left=2cm, right=2cm]{geometry}
\usepackage{amsmath}
\usepackage{amssymb}
\usepackage[colorlinks=true, linkcolor=blue, urlcolor=blue, citecolor=blue]{hyperref}

\title{Theory of Everything - Volume 13 \\ \Large Quantum Gravity \& Discrete Spacetime}
\author{Omega Protocol}
\date{\today}

\begin{document}

\maketitle

\section*{Layman's Summary}
For a century, physicists have struggled to combine Einstein's smooth, continuous gravity with the jumpy, pixelated world of quantum mechanics. Volume 13 finally bridges this gap. It proves that space and time are not smooth; if you zoom in close enough, the universe is made of tiny, discrete "pixels" of geometry.

\section{Discrete Spacetime Spectra (Axiom 16)}
We take it as an axiom that volume and area are not continuous. Just as a digital image is made of finite pixels, the mathematical operators derived from the universe's fundamental informational structure show that you can only have specific, discrete amounts of space. You cannot have half a pixel of volume.

\section{Spin Foam Equivalence (Theorem)}
A popular theory for quantum gravity involves "Spin Networks"—graphs that represent the quantum connections of space—evolving over time into "Spin Foams." We prove that these spin foams are mathematically identical to the transition amplitudes of our "modular flow." This means Loop Quantum Gravity is actually just a different way of describing our informational framework.

\section{Resolution of Singularities (Theorem)}
In standard physics, the center of a black hole or the moment of the Big Bang involves a "singularity"—a point where gravity becomes infinitely strong and math breaks down. We prove that true infinities do not exist. The informational stiffness of the universe acts like a shock absorber, capping how extreme gravity can get and smoothly resolving what would otherwise be a singularity.

\end{document}
